\documentclass{leveldoc}
\DocUses{foundations,slingshot,first_coil}
\DocFingerprint{10b28fda106b1bb1}

\begin{document}

\section{What the level asks}\label{sec:asks}

The game's text:
\leveldescription

Six positive ions leave a source at $(\val{xZero}, \val{yZero})$ with the small kinetic energy
$T_0 = \val{Tzero}$: three masses, $1$, $2$ and $4$, each at $\pm\val{angDeg}$ degrees from
$+y$. The board lies inside a large circular coil whose field, perpendicular to the board,
bends them round. Each ion must enter the detector of its mass, all on the line $y = 5$ to the
right of the source (the table in \S\ref{sec:hand}); there are no acceptance windows, kinetic
windows or gates. The player places up to three charges from the listed magnitudes in the
region $\val{regX0} \le x \le \val{regX1}$, $\val{regY0} \le y \le \val{regY1}$, below the
source. A flight ends when the ion enters its detector, touches a charge or the coil's wire,
leaves the world, or reaches $t_{\max}$ (\xsecref{foundations}{sec:board}).

The reference places one positive charge $Q = \val{QOne}$ at $(\val{xOne}, \val{yOne})$, just
below the source: its repulsion accelerates the ions upwards, and the coil's field bends them
back down after half a turn onto the detectors, the heavier ions on wider circles. This is
Dempster's spectrometer: an accelerating field, then a magnetic half-turn whose diameter
grows as the square root of the mass, with the two launch angles re-converging after half a
turn. This document estimates the diameters in a uniform field and integrates the six flights
numerically.

\section{Setup}\label{sec:setup}

The level file, field by field (the appendix reproduces it):
\setuptable

In symbols: every ion has charge $q = \val{q}$ and radius $0$; $c = \val{c}$. The coil is a
circle of radius $a = \val{aCoil}$ about $(\val{xCoil}, \val{yCoil})$ with the strength
$\kappa = \val{kappa}$ (\physref{2.2}), whose field at its centre is $2\pi\kappa/a =
\val{Bcentre}$ (\xsecref{first_coil}{sec:loop}). The reference charge has radius $R = \val{R}$,
$\val{rLaunch}$ below the source. The world is $-1 \le x \le \val{xWorld}$,
$-1 \le y \le \val{yWorld}$.

\begin{figure}[htbp]
\centering
\begin{tikzpicture}
\begin{axis}[axis equal image, width=0.92\linewidth, xmin=-1.6, xmax=33.6, ymin=-1.6, ymax=23.6,
  xtick={0,5,...,30}, ytick={0,5,...,20}, tick label style={font=\small},
  xlabel={$x$}, ylabel={$y$}, label style={font=\small}]
  \draw[black!8] (0,0) grid (32,22);
  \draw[black!45, dashed] (-1,-1) rectangle (\rawval{xWorld},\rawval{yWorld});
  \draw[black!70, thick] (\rawval{xCoil},\rawval{yCoil}) circle[radius=\rawval{aCoil}];
  \draw[black!30, dotted] (\rawval{regX0},\rawval{regY0}) rectangle (\rawval{regX1},\rawval{regY1});
  \fill[green!50!black, fill opacity=0.18] (13,2) rectangle (14,5);
  \draw[green!40!black] (13,2) rectangle (14,5);
  \fill[green!50!black, fill opacity=0.18] (17,2) rectangle (18,5);
  \draw[green!40!black] (17,2) rectangle (18,5);
  \fill[green!50!black, fill opacity=0.18] (22,2) rectangle (24,5);
  \draw[green!40!black] (22,2) rectangle (24,5);
  \addplot[blue!70!black] table {patha.dat};
  \addplot[blue!70!black, dashed] table {pathb.dat};
  \addplot[red!70!black] table {pathc.dat};
  \addplot[red!70!black, dashed] table {pathd.dat};
  \addplot[orange!80!black] table {pathe.dat};
  \addplot[orange!80!black, dashed] table {pathf.dat};
  \fill[red!75!black] (\rawval{xOne},\rawval{yOne}) circle[radius=\rawval{R}];
  \fill (\rawval{xZero},\rawval{yZero}) circle[radius=0.15];
\end{axis}
\end{tikzpicture}
\caption{The coil (black circle), the reference charge (red disc) below the source, and the six
half-turns onto the detectors (green): mass 1 (blue), 2 (red) and 4 (orange), each at
$+\val{angDeg}^\circ$ (solid) and $-\val{angDeg}^\circ$ (dashed). The detectors' positions are
the level file's (the table of \S\ref{sec:hand}).}
\label{fig:board}
\end{figure}

\section{Principles and model}\label{sec:model}

\begin{principles}
  \item The Lorentz force in the relativistic equation of motion \xeqref{foundations}{eq:eom},
    the point charge's field \xeqref{foundations}{eq:coulomb}, the energy
    \xeqref{foundations}{eq:energy}; the magnetic force does no work
    (\xsecref{first_coil}{sec:nowork}).
  \item The coil's field in its plane, \xeqref{first_coil}{eq:loop-field}, with $K$ and $E$ by
    \xeqref{first_coil}{eq:agm}; the circle of \xeqref{first_coil}{eq:larmor-radius}.
  \item The board's rules; the classical Runge--Kutta method, cubic Hermite interpolation and
    the error estimate by doubling the step (\xsecref{foundations}{sec:board},
    \xsecref{foundations}{sec:numerics}), as in \xsecref{slingshot}{sec:method}.
\end{principles}

The assumptions are First coil's (\xsecref{first_coil}{sec:model}): the coil and the charge are
held in place, the coil is a thin circle, the ions are points, and their radiation is not taken
from them.

\section{Derivation}\label{sec:derivation}

\subsection{The acceleration}\label{sec:accel}

The coil's field does no work; the charge does. By the energy \xeqref{foundations}{eq:energy}
an ion's kinetic energy is $T = T_0 + qQ\,(1/r_\mathrm{launch} - 1/r)$, with $r$ its distance from
the charge. At the source $qQ/r_\mathrm{launch} = \val{VLaunch}$, twenty times $T_0$: the ions
leave almost at rest and take their energy from the charge, all of them about the same, whatever
their mass. Mass 1 has $T = \val{TtopOne}$ at the top of its half-turn ($y = \val{topY}$).

\subsection{Diameters and focusing in a uniform field}\label{sec:diameter}

In a uniform field $B$ an ion of momentum $p$ turns on a circle of radius $p/\abs{qB}$
\xeqref{first_coil}{eq:larmor-radius}, and after half a turn it is a diameter away from where it
started:
\begin{equation}\label{eq:diameter}
  D = \frac{2p}{\abs{qB}} \approx \frac{2\sqrt{2mT}}{\abs{qB}} \propto \sqrt{m}
\end{equation}
at equal energy ($p \approx \sqrt{2mT}$ at these speeds, $T \ll mc^2$): the masses $1$, $2$ and
$4$ land at distances in the ratio $1 : \sqrt2 : 2$. An ion launched at the angle $\alpha$ from
the central direction lands at $D\cos\alpha$, to second order $D - D\alpha^2/2$ for either sign of
$\alpha$: the two angles land together, $180^\circ$ focusing. Against the central ray they fall
short by $r\alpha^2$, here $\val{focusOne}$ for mass 1 ($\alpha = \val{alphaRad}$).

The coil's field is not uniform (\xsecref{first_coil}{sec:loop}): it grows towards the wire,
and the charge's field still pushes along the way. With $p$ and $B_z$ taken at the top of each
mass's half-turn from the march below (the $+\val{angDeg}^\circ$ ions):
\begin{center}\small\val{estimates}\end{center}
For mass 1 the estimate lands within a few per cent; the heavier ions' wider circles reach the
stronger field near the wire and land shorter. In a uniform field the two angles would land
together; here their landing points differ by $\val{spreadOne}$ (mass 1), $\val{spreadTwo}$
(mass 2) and $\val{spreadFour}$ (mass 4), through the field's gradient and the charge, but by
less than the detectors' widths.

\subsection{The method}\label{sec:method}

The six flights are integrated by \xeqref{foundations}{eq:rk4} at $h = \val{H}$ in the coil's
field and the charge's, as in \xsecref{first_coil}{sec:method}; the energy is the check. The
largest step-doubling estimate of an end state is $\val{errMax}$, the largest energy drift
$\val{wDevMax}$.

\section{The reference solution by hand}\label{sec:hand}

\begin{procedure}
  \item\label{st:launch} \emph{The launches.} $T_0 = \val{Tzero}$ for every ion, at
    $\pm\val{angDeg}^\circ$ from $+y$; each mass's momentum from
    \xeqref{foundations}{eq:launch}.
  \item\label{st:march} \emph{The march.} Advance each ion by \xeqref{foundations}{eq:rk4} in
    the coil's field \xeqref{first_coil}{eq:loop-field} and the charge's, and locate its entry:
    \begin{center}\small\val{results}\end{center}
    Every ion enters its detector through its top face, $y = 5$, inside the detector's span.
\end{procedure}

\begin{result}[The reference flights]
All six ions enter the detectors of their masses at the points of the table. The level is
solved.
\end{result}

\section{Comparison with the game}\label{sec:comparison}

\comparisontable

All differences are a few $10^{-11}$ or less, the game's verified accuracy (Foundations §6);
our step estimates are a few $10^{-13}$. The tolerance, $10^{-9}$, is ten times the larger of
the two computations' estimates (or the game's own preview--verification gap, where larger),
fixed before comparing.

\section{Other solutions}\label{sec:others}

On the empty board, the second placement in \code{placements.json}, the ions leave the source
almost at rest and circle slowly in the coil's field, drifting along its gradient; the light
ones still drift into their detectors after long times, but the heavy ones do not before
$t_{\max}$ (ion 5 ends as \val{bareOutFive}): the comparison table gives each outcome. A charge
must give the ions their energy, and one below the source, as in the reference, gives all of
them nearly the same, which \eqref{eq:diameter} needs. A map of solutions would take the game's
verdicts, which this document does not reproduce.

\section{What the game computes}\label{sec:engine}

The game evaluates the coil's field by its closed form (\physref{2.2}) and integrates the
equations of motion with its adaptive eighth-order Runge--Kutta method (\physref{5},
\physref{6}, \physref{7}); the differences in the comparison are the two integrations' errors.

\section{Further reading}\label{sec:reading}

Motion in a magnetic field: \Jackson{§12.2}; the coil's field and cyclotron motion:
\levelref{first_coil}. The instrument is A. J. Dempster's mass spectrometer of 1918, the
first with $180^\circ$ magnetic focusing.

\levelfileappendix

\end{document}
